\documentclass[11pt]{article}

\usepackage[margin=1in]{geometry}
\usepackage{amsmath, amssymb}
\usepackage{xcolor}
\usepackage{listings}
\usepackage{hyperref}
\usepackage{array}
\usepackage{booktabs}

\hypersetup{
    colorlinks=true,
    linkcolor=blue!60!black,
    urlcolor=blue!60!black,
    citecolor=blue!60!black
}

\lstdefinestyle{bugpython}{
  language=Python,
  basicstyle=\ttfamily\small,
  keywordstyle=\color{blue!60!black},
  stringstyle=\color{green!40!black},
  commentstyle=\color{gray!70!black},
  showstringspaces=false,
  breaklines=true,
  frame=single,
  rulecolor=\color{gray!30},
  columns=fullflexible,
  keepspaces=true
}

\title{SymPy Bug Report}
\author{}
\date{}

\begin{document}

\maketitle

\section{Bug 27: solveset returns a false acosh solution outside the principal range}

\subsection{Minimal Reproducer}

\medskip

\begin{lstlisting}[style=bugpython]
from sympy import Eq, I, N, S, acosh, pi, solveset, symbols

x = symbols("x")
eq = Eq(acosh(x), 2*pi*I)
sol = solveset(eq, x, domain=S.Complexes)
print("solution =", sol)
for s in sol:
    print("residual at", s, "=", N(acosh(s) - 2*pi*I, 80))
\end{lstlisting}

\subsection{Observed Behavior}

\medskip

\begin{lstlisting}[style=bugpython]
solution = {1}
residual at 1 = -6.283185307179586476925286766559005768394338798750211641949889184615632812572418*I
\end{lstlisting}

SymPy returns the finite solution set \(\{1\}\).  Direct substitution shows that
this returned point does not satisfy the original equation: the residual is
\(-2\pi i\), displayed above numerically to 80 digits.

\subsection{Expected Behavior}

The correct result is \(\varnothing\).  The principal value satisfies
\(\operatorname{acosh}(1)=0\), not \(2\pi i\), and there is no complex number
\(x\) whose principal inverse hyperbolic cosine is \(2\pi i\).

\subsection{Mathematical Explanation}

The error comes from treating \(\cosh\) as a global inverse of the principal
function \(\operatorname{acosh}\).  Although
\[
  \cosh(2\pi i)=1,
\]
this only proves that \(1\) solves the composed equation
\[
  \cosh(\operatorname{acosh}(x)) = \cosh(2\pi i).
\]
It does not prove that \(1\) solves the original principal-branch equation
\[
  \operatorname{acosh}(x)=2\pi i.
\]
For the returned candidate,
\[
  \operatorname{acosh}(1)-2\pi i = -2\pi i \ne 0.
\]
Thus the result is not an alternative branch representation or a harmless
simplification; it is a false solution introduced by discarding the range
condition of the principal inverse hyperbolic cosine.

\subsection{Source-Level Diagnosis}

The diagnosis identifies the relevant solver logic in
{\footnotesize\nolinkurl{sympy/solvers/solveset.py}}.  The key function is
\texttt{\_invert\_complex} around lines 550--557, with the call reaching
\texttt{\_solveset} around line 1313.  The traced call path rewrites the input
as \(\operatorname{acosh}(x)-2\pi i=0\), removes the additive
\(-2\pi i\), and then inverts
\(\operatorname{acosh}(x)=2\pi i\) through the generic complex inversion path.
That path asks \(\operatorname{acosh}(x)\) for its inverse.  The inverse
method for \(\operatorname{acosh}\), located in the file
{\footnotesize\nolinkurl{sympy/functions/elementary/hyperbolic.py}}, returns
\(\cosh\).

The problematic rule is the generic mapping of the target set through the
function returned by \texttt{inverse()}:

\begin{lstlisting}[style=bugpython]
if hasattr(f, 'inverse') and f.inverse() is not None:
    return x, imageset(Lambda(n, f.inverse()(n)), g_ys)
\end{lstlisting}

For this equation, that rule maps the target set \(\{2\pi i\}\) through
\(\cosh\), producing \(\{\cosh(2\pi i)\}=\{1\}\).  The diagnosis notes that
this treats \(\cosh\) as though it were a globally valid inverse for
principal \(\operatorname{acosh}\).  Since \(\operatorname{acosh}\) has a
principal range, the inversion is only sound when the right-hand side is in
that range.  Here it is not, so the generic inversion solves the composed
periodic equation instead of the original equation.

A plausible fix direction supported by the diagnosis is to exclude inverse
hyperbolic functions from this branch-blind generic inverse rule, or to add
range conditions for the target values before mapping through the elementary
inverse.  The diagnosis also notes that a residual check on finite candidates
would prevent this particular false solution from being returned.

\subsection{Additional Failing Instances}

The independent verification covers the representative target \(2\pi i\) and
five additional targets of the same form.  In each case SymPy maps the target
through \(\cosh\), returning \(1\) or \(-1\), but the principal value of
\(\operatorname{acosh}\) at the returned candidate is \(0\) or \(\pi i\), not
the requested target.  The expected behavior is therefore the same across the
family: the listed false candidate must not be reported as a solution.

\begin{center}
\footnotesize
\renewcommand{\arraystretch}{1.3}
\begin{tabular}{@{}>{\raggedright\arraybackslash}p{0.38\linewidth}>{\raggedright\arraybackslash}p{0.25\linewidth}>{\raggedright\arraybackslash}p{0.27\linewidth}@{}}
\toprule
\textbf{Input / Expression} & \textbf{SymPy behavior} & \textbf{Expected behavior} \\
\midrule
\(\operatorname{solveset}(\operatorname{acosh}(x)=2\pi i)\) & returns a set containing \(1\) & \(1\) is not a solution; \(\operatorname{acosh}(1)=0\) \\
\(\operatorname{solveset}(\operatorname{acosh}(x)=-2\pi i)\) & returns a set containing \(1\) & \(1\) is not a solution; \(\operatorname{acosh}(1)=0\) \\
\(\operatorname{solveset}(\operatorname{acosh}(x)=3\pi i)\) & returns a set containing \(-1\) & \(-1\) is not a solution; \(\operatorname{acosh}(-1)=\pi i\) \\
\(\operatorname{solveset}(\operatorname{acosh}(x)=-3\pi i)\) & returns a set containing \(-1\) & \(-1\) is not a solution; \(\operatorname{acosh}(-1)=\pi i\) \\
\(\operatorname{solveset}(\operatorname{acosh}(x)=4\pi i)\) & returns a set containing \(1\) & \(1\) is not a solution; \(\operatorname{acosh}(1)=0\) \\
\(\operatorname{solveset}(\operatorname{acosh}(x)=-4\pi i)\) & returns a set containing \(1\) & \(1\) is not a solution; \(\operatorname{acosh}(1)=0\) \\
\bottomrule
\end{tabular}
\end{center}

\subsection{Related Issues Assessment}

The best-effort deduplication check found public material on the relevant
principal-branch mathematics for inverse trigonometric and inverse hyperbolic
functions, but did not identify a public SymPy issue or fix for the exact
\(\operatorname{acosh}(x)=2\pi i\) false-solution reproducer.  The verdict is
therefore a new instance of a known failure mode: the broader range-condition
issue is recognizable, but this concrete reproducible behavior is previously
unreported in the available evidence and remains individually actionable as a
SymPy issue and regression test.

\subsection{Proposed SymPy Regression Test}

\medskip

\begin{lstlisting}[style=bugpython]
import pytest
from sympy import Eq, I, S, acosh, pi, simplify, solveset, symbols

x = symbols("x")

CASES = [
    (2*pi*I, 1),
    (-2*pi*I, 1),
    (3*pi*I, -1),
    (-3*pi*I, -1),
    (4*pi*I, 1),
    (-4*pi*I, 1),
]


@pytest.mark.parametrize("target,bad", CASES)
def test_solveset_acosh_rejects_values_outside_principal_range(target, bad):
    sol = solveset(Eq(acosh(x), target), x, domain=S.Complexes)

    assert sol.contains(bad) is not S.true
    assert simplify(acosh(bad) - target) != 0
\end{lstlisting}

\end{document}
